\chapter*{不定积分每日一题}

\refer{《积分的方法与技巧》金玉明\ 顾新身\ 毛瑞庭\ /\ 编著}

\exercise{求积分$\frac{\sin{x}}{a\cos{x}+b\sin{x}}$}%在花括号内输入题目

\teacher{张晓鹏}%在花括号内输入你的名字

\postdate{2023年4月10日}%在花括号内输入你想要发布那一天的日期

\result{\frac{1}{a^2+b^2}\qty(bx-a\ln{\abs{a\cos{x}+b\sin{x}}})+C}%在花括号内输入本题结果

\analysis{
    对于这个题目，很明显符合我们昨天万能替换法，感兴趣的同学可以试试，不过实际上属于“试试就逝世”，计算量太大了，而对于这种比较“规则”的积分，我们可以给他构造出另外一个长得很像的积分，放在一起配对积分，从而化简计算.

    比如这题，我们可以构造积分：
    \formula{I_1=\int{\frac{\sin{x}}{a\cos{x}+b\sin{x}}},\ I_2=\int{\frac{\cos{}}{a\cos{x}+b\sin{x}}}\d x}
    将这两个积分组合起来，我们可以得到
   \formula{
    \begin{aligned}
        bI_1+aI_2&=\int{\frac{b\sin{x}+a\cos{x}}{a\cos{x}+b\sin{x}}}=\int{\d x}=x+C\\
        -aI_2+bI_2&=\int{\frac{-a\sin{x}+b\cos{x}}{a\cos{x}+b\sin{x}}\d x}\\
        &=\int{\frac{\d\qty(a\cos{x}+b\sin{x})}{a\cos{x}+b\sin{x}}}\\
        &=\ln{\abs{a\cos{x}+b\sin{x}}}+C
    \end{aligned}
   }
   那么剩下的过程相当于解方程，从两式中消去$I_2$，即可得到原积分结果：
   \formula{\qty(a^2+b^2)I_1=bx-a\ln{\abs{a\cos{x}+b\sin{x}}}+C}
   \formula{\Rightarrow I_1=\frac{bx-a\ln{\abs{a\cos{x}+b\sin{x}}}}{a^2+b^2}+C}
   而对应的另外一个积分$I_2$也可以得到
   \formula{\Rightarrow I_2=\frac{ax+b\ln{\abs{a\cos{x}+b\sin{x}}}}{a^2+b^2}+C}
   而这就是组合积分法，组合积分法的使用场景主要是三角函数、指数函数、双曲函数 这些互导或自导函数中的一些特定形式，而对应的特殊形式需要靠大家积累了，一句两句的确实难以清楚也无法完全表达全部.
}%在花括号内输入分析